% 本文件由 LaTeX 排版 Agent 根据有效 translation.json 自动生成。
% author=Justin Martyr; work_id=0127; title=圣犹斯定第二护教书

\chapter{圣犹斯定第二护教书}

% structure: p0001 source=h1 target=omitted reason=page-title-already-rendered-by-parent

致罗马元老院\par

% structure: p0003 source=h2 target=section reason=relative-heading-rank; base=h2

\section{第一章 引言}

罗马人啊，最近在你们城中，在\Person{乌尔比古斯}治下所发生的事，以及同样在各处由地方官不理智地行出的事，迫使我为了你们的缘故作此篇陈辞。你们是同样有情欲的人，亦是弟兄，虽然你们不知道，并且因你们所尊崇的显位而夸耀，不肯承认。因为各处，凡被父亲、邻居、儿女、朋友、兄弟、丈夫或妻子，因过失、顽固、贪图享乐、难以劝勉归正而加以纠正的人（除了那些确信不义与放纵者将在永火中受罚，而有德者与像基督一样生活的人，将无痛苦地与天主同住——我指的是已成为基督徒的人），以及憎恨我们的邪恶魔鬼，使这样的人臣服于自己，以审判官的身份役使他们，作为被恶灵驱动的统治者，煽动他们置我们于死地。但为了使你们清楚在\Person{乌尔比古斯}治下所发生的一切原因，我将叙述所发生的事。\par

% structure: p0005 source=h2 target=section reason=relative-heading-rank; base=h2

\section{第二章 \Person{乌尔比古斯}判处基督徒死刑}

某妇女曾与一放纵的丈夫同居；她本人先前也曾放纵。但当她知道基督的教导后，便变得清醒，并努力说服丈夫也节制，引用基督的教导，向他保证那些不节制、不符合正确理性而生活的人将在永火中受罚。但他继续同样的放荡，因自己的行为使妻子疏远了他。因她认为继续与一个想尽办法放纵享乐、违反自然律并违背正义的丈夫同住是邪恶的，便希望与他离婚。当她被朋友们劝说继续与他同住，认为丈夫可能某时会有改过的希望时，她勉强了自己的情感，仍与他同居。但当丈夫去了\Place{亚历山大}，并被传行为比以前更糟时，她——为避免因继续与他保持婚姻关系、共享餐桌与床榻而成为他邪恶与不敬的共犯——给了他你们所谓的离婚书，并与他分离。但这高贵的丈夫——他本应欢喜，因她从前毫无顾忌地与仆人和雇工所做的那些事，当她以醉酒和各样恶行为乐时，如今都已放弃，并希望他也放弃同样的事——当她离开他而非出于他的意愿时，他却控告她，声称她是基督徒。她向皇帝您呈递了一份文书，请求先允许她处理事务，待事务安排妥当后，再就控告进行辩护。您准允了此请求。而她从前的丈夫，因不能再控告她，便将攻击转向一人，即\Person{托勒密乌斯}，此人曾是她基督道理的教师，被\Person{乌尔比古斯}惩罚。他以下列方式行事：他说服一位百夫长——此百夫长曾将\Person{托勒密乌斯}下狱，且与他友善——将\Person{托勒密乌斯}提来，只盘问这一点：他是否是基督徒？\Person{托勒密乌斯}是爱真理之人，并无欺诈或虚伪的性情；当他宣认自己是基督徒时，便被百夫长捆绑，在监狱中长期受刑。最后，当此人被带到\Person{乌尔比古斯}面前时，他只被问了这一个问题：他是否是基督徒？他再次意识到自己的责任，并因基督的教导而认识到其中之高贵，便宣认自己在神圣美德中的门徒身份。因为否认某事的人，要么因谴责其事本身而否认，要么因意识到自己不配或与之疏远而退缩于宣认；这两种情况都不属于真基督徒。当\Person{乌尔比古斯}下令将他带去受刑时，一位同样也是基督徒的\Person{路求}，看到如此不合理的判决，便对\Person{乌尔比古斯}说：\Quote{这判决的依据是什么？你为什么惩罚这个人？他不是奸夫、淫乱者、杀人犯、窃贼、强盗，也没有被判定任何罪行，只承认被称为基督徒的名字。\Person{乌尔比古斯}啊，你的这个判决与皇帝庇护、哲学家恺撒之子，以及神圣元老院都不相称。}他对\Person{路求}的答复只说了这句：\Quote{我看你也是这样的人。}当\Person{路求}回答：\Quote{我当然也是。}时，他又下令将他也带走。他表达了感谢，知道自己得以脱离如此邪恶的统治者，正前往天上的圣父与君王那里。还有第三个人上前，也被判受刑。\par

% structure: p0007 source=h2 target=section reason=relative-heading-rank; base=h2

\section{第三章 \Person{犹斯定}指控\Person{克雷斯森斯}对基督徒的无知偏见}

因此，我也预料会被我所提到的那些人，或者被\Person{克雷森斯}——那个喜欢虚张声势和夸口的人——设计陷害并被钉在火刑柱上。因为这个人不配称为哲学家，他公开在不明真相的事情上作证反对我们，说基督徒是无神论者和不敬虔的人，这样做是为了赢得受蒙蔽群众的欢心并取悦他们。如果他未曾阅读基督的教导就攻击我们，那他就是彻底堕落，远胜于无知之辈，后者往往避免讨论或作伪证于他们所不明白的事。或者，如果他读过这些教导却不明白其中的威严，或者明白了却这样做以免被怀疑是（这样一个）基督徒，那他更是卑鄙无耻、彻底堕落，被卑劣和不合理的意见与恐惧所征服。因为我想让你们知道，我曾就这个问题向他提出过一些问题，并质问他，结果使我深信他实际上什么都不知道。为了证明我说的是实话，如果这些讨论还没有报告给你们，我准备在你们面前重新进行。这将是一件配得上君主的行为。但如果我的问题和他的回答已经为你们所知，那你们已经明白他对我们的事一无所知；或者，如果他了解这些事，却因害怕那些可能听见他说话的人而不敢直言，那他就如苏格拉底一样，正如我先前所说，证明自己不是哲学家，而是一个固执己见的人；至少他不尊重那句苏格拉底式的、最值得钦佩的话语：\Quote{但一个人绝不可在任何事情上被尊崇于真理之前。}然而，对于一个以漠然为目标的犬儒者来说，除了漠然之外，不可能认识任何善。\par

% structure: p0009 source=h2 target=section reason=relative-heading-rank; base=h2

\section{第四章 为什么基督徒不自杀}

但免得有人对我们说：\Quote{那么你们所有人都去自杀吧，现在就回到天主那里，不要来麻烦我们。}我要告诉你们为什么我们不这样做，为什么我们在受审时无畏地宣认。我们被告知，天主并非毫无目的地创造世界，而是为了人类；我们之前也说过，祂喜悦那些效法祂属性的人，不喜悦那些在言语或行为上拥抱无价值之事的人。那么，如果我们都自杀了，我们就会成为，就我们而言，无人出生、无人受教于神圣教义、甚至人类种族不复存在的原因；而且，如果我们如此行，我们自身就在违反天主的旨意。但当受审时，我们不否认，因为我们不觉得自己有恶，反而认为在一切事上不说真话是不敬虔的，我们也知道这蒙天主喜悦，并且我们现在非常渴望将你们从一种不公正的偏见中解救出来。\par

% structure: p0011 source=h2 target=section reason=relative-heading-rank; base=h2

\section{第五章 天使如何堕落}

但如果有人产生这样的想法：既然我们承认天主是我们的帮助者，我们就不应该如我们所说被恶人压迫和迫害——这一点我也要解答。天主创造了整个世界，将地上万物置于人的管辖之下，安排了天上的元素以促进果实的生长和季节的更替，并制定了这条神圣的法律（因为祂显然也为人类创造了这些事），祂将人类和天下一切事物的管理交给了祂所设立的天使。但天使们违背了这一任命，被对女人的爱所俘获，生下了被称为魔鬼的子女；此外，他们后来通过魔法著作、恐惧和所施加的惩罚，以及教导人献祭、焚香和奠酒（这些东西在他们被情欲的冲动所奴役后便需要），部分地使人类臣服于自己；并在人类中播下了凶杀、战争、奸淫、放纵的行为和一切邪恶。因此，诗人和神话学家们，不知道是这些天使和由他们所生的魔鬼对男人、女人、城市和民族做了他们所叙述的那些事，便将它们归于天主本身，以及被认为是祂亲生儿女的那些，以及被称为祂兄弟的尼普顿和普路托的子女，再以及这些子孙的子女。因为每个天使给自己和他的子女起的名字，他们就以那名字称呼他们。\par

% structure: p0013 source=h2 target=section reason=relative-heading-rank; base=h2

\section{第六章 天主和基督的名字、意义与能力}

对于万有的圣父，祂是非受生的，并没有一个名字。因为无论祂被称作什么名字，那给祂起名的人都是祂的先行者。但\Quote{圣父}、\Quote{天主}、\Quote{创造者}、\Quote{主}和\Quote{主宰}这些词都不是名字，而是源自祂的善行和功用的称谓。而祂的圣子，那唯独真正被称为圣子的，那圣言，祂本与祂同在，并在创造万有之先受生，当祂起初借着祂创造并安排一切时，祂被称为基督，是指祂受傅油和天主借着祂安排一切；这名字本身也含有一个未知的意义；同样，\Quote{天主}这个称谓也不是名字，而是植根于人类本性中对一个难以解释之物的观念。而\Quote{耶稣}，作为人和救主的名字，也有其意义。因为祂也成了人，如我们先前所说，按照天主圣父的旨意受孕，是为了相信的人，并为了毁灭魔鬼。现在你们可以从你们所观察到的事中学到这一点。因为全世界无数附魔者，在你们的城市里，我们许多基督徒奉那在本丢·比拉多手下被钉十字架的耶稣基督的名驱魔，已经治愈并仍在治愈，使附着于人的魔鬼无力并被逐出，尽管所有其他驱魔者以及使用咒语和药物的人都不能治愈他们。\par

% structure: p0015 source=h2 target=section reason=relative-heading-rank; base=h2

\section{第七章 世界因基督徒而存留。人的责任}

因此，天主要使整个世界陷入混乱和毁灭——藉此，恶天使、魔鬼和人都将不复存在——之所以延迟此事，是因为基督徒的种子，他们知道自己是在自然界中保存的原因。因为若不是这样，你们就不可能做出这些事，也不可能被邪灵驱使；然而审判的火必要降下，彻底消溶一切，正如从前洪水只留下了我们称为\Person{诺厄}、你们称为\Person{丢卡利翁}的那一人及其全家，从他们又生出如此众多的人，有的邪恶，有的良善。因为我们如此说，将有那焚毁，但并非如斯多葛派所主张的万物互相转化的理论——那似乎十分可鄙。但我们也不主张，人的作为或所受的遭遇是出于命运，而是说每人凭自由选择行善或犯罪；并且是由于魔鬼的影响，像\Person{苏格拉底}等认真的人遭受迫害和监禁，而\Person{萨达那帕鲁斯}、\Person{伊壁鸠鲁}等之辈却似乎因富足和荣耀而蒙福。斯多葛派没有看到这一点，便主张一切事情都按命运的必然性发生。但由于天主在起初使天使和人类种族都有自由意志，他们将公正地因自己所犯的罪而在永火中受罚。这是所有受造之物的本性：能够有恶与德。因为若非有转向两者的能力，就没有一个可受称赞。这一点也由各处那些依照正理立法和从事哲学的人所表明，他们规定要做某些事，避免另一些事。甚至斯多葛派的哲学家，在其道德学说中，也坚定地尊崇同样的事，因此显然他们关于原理和无形之物的说法并不很成功。因为若他们说人的行为是命运所致，那么他们就会主张：天主不外乎是那些永远变化、改变并消溶为相同事物之物，于是他们似乎只理解了可朽坏之物，并且认为天主本身在每一个恶事中部分或全部地出现；或者主张恶与德都毫无意义——这是与一切健全的思想、理性和感觉相悖的。\par

% structure: p0017 source=h2 target=section reason=relative-heading-rank; base=h2

\section{第八章 凡圣言寓居者皆被憎恨}

而斯多葛学派的人——就其道德教训而言，他们是可钦佩的，如同诗人在某些方面也一样，因为理性（圣言）的种子植根于每一人类种族——我们知道，他们被憎恨并被处死，例如\Person{赫拉克利特}，以及我们时代的\Person{穆索尼乌斯}等人。因为正如我们所说，魔鬼常使那些按照理性与认真生活、远离邪恶的人被憎恨。这并不奇怪；如果魔鬼被证实使那些不仅按所散布的圣言之一部分生活，而是因认识并默观整个圣言（即基督）而生活的人更加被憎恨。他们被封在永火中，将遭受公义的惩罚。因为倘若他们如今已藉耶稣基督的名被人打倒，这便是一种预兆，表明他们自己和服事他们的人将要受永火的刑罚。因为众先知都如此预言，我们的导师\Person{耶稣}也如此教导。\par

% structure: p0019 source=h2 target=section reason=relative-heading-rank; base=h2

\section{第九章 永罚非虚张声势}

为使没有人像那些被视为哲学家的人所说的那样，声称我们关于恶人受永火惩罚的断言是夸夸其谈和吓唬人，并说我们希望人因恐惧而度道德生活，并非因为此种生活是美善愉快的；我将对此作简要答复：倘若并非如此，则天主不存在；或者若祂存在，祂并不关心人类，那么德与恶都毫无意义，并且如我们先前所说，立法者惩罚违反善法令的人是不义的。但因这些人并非不义，他们的圣父借着圣言教导他们做同样的事，与他们一致的人并非不义。若有人反对说，人的法律各不相同，并说在有些地方，一件事被认为善，另一件为恶；而在别处，前者认为恶的却被视为善，前者认为善的被视为恶，就让他听我们对此怎么说。我们知道，恶天使依照自己的邪恶制定了法律，与他们相似的人喜爱这些法律；而当正理降临时，祂证明并非所有意见和所有学说都是善的，而是有些是恶的，有些是善的。因此，我将向这样的人宣布同样的和类似的事，若有需要，将更详细地讨论。但眼下我回到本题。\par

% structure: p0021 source=h2 target=section reason=relative-heading-rank; base=h2

\section{第十章 基督与苏格拉底之比较}

我们的教义，因而，显得比一切人的教导更伟大；因为基督，为了我们而显现，成了整个有理性的存在，即身体、理智和灵魂。因为我，无论立法者或哲学家所说的好话，他们是通过发现和观察\OriginalTerm{圣言}{Word}的某一部分而精心阐发的。但由于他们，我并不知道整个\OriginalTerm{圣言}{Word}，即基督，他们常常自相矛盾。而那些按照人的出生比基督更古老的人，当他们试图通过理智思考和证明事物时，被带到法庭前，被当作不敬神的人和好事者。而\Person{苏格拉底}，在这方面比他们所有人都更热心，被指控与我们所受的完全相同的罪。因为他们说他引入新神，并且不把国家所承认的那些当作神。但他将\Person{荷马}和其他诗人逐出国境，教导人们拒绝那邪恶的魔鬼以及那些做诗人所描述之事的人；并劝勉他们通过理性的探究来认识那位对他们而言不认识的天主，说：\Quote{要找到万有的\OriginalTerm{圣父和造物主}{Father and Maker}，既不容易，即使找到了，向所有人宣告祂也是不稳妥的。}但这些事我们的基督藉着自己的能力成就了。因为没有人信赖\Person{苏格拉底}到愿意为这教义而死的地步，但在基督里，祂甚至在苏格拉底那里也被部分地认识（因为祂曾是并永远是那在每个人内的圣言，祂曾通过先知预言将要发生的事，并且当祂成为与我们相似的情欲时，亲自教导了这些事），不但哲学家和学者相信了，就连工匠和完全没有教育的人也相信了，他们藐视荣耀、恐惧和死亡；因为祂是那不可言说的圣父的能力，而非单纯的人类理性的工具。\par

% structure: p0023 source=h2 target=section reason=relative-heading-rank; base=h2

\section{第十一章 基督徒如何看待死亡}

但我们也不应被处死，恶人和魔鬼也不会比我们更有力，要不是死亡是每个生者所欠的债。因此，当我们偿还这笔债时，我们心怀感恩。我们认为在此处为了\Person{克雷森斯}和他的那些如他一样狂乱的人，讲述\Person{色诺芬}所记载的事是合宜而有益的。\Person{色诺芬}说，\Person{赫拉克勒斯}来到一个三叉路口，发现\OriginalTerm{德行}{Virtue}和\OriginalTerm{恶行}{Vice}以妇女形象出现在他面前：\OriginalTerm{恶行}{Vice}穿着奢华的衣饰，带着诱人的表情，因这样的装饰而显得容光焕发，她的眼睛流露出迅速融化的温柔，她对\Person{赫拉克勒斯}说，如果跟随她，她将使他终身在享乐中度过，并装饰以最优雅的饰物，正如当时她自己身上所穿戴的那样；而\OriginalTerm{德行}{Virtue}则形容枯槁，衣着褴褛，说道：但你若服从我，你将不以那逝去和朽坏的装饰或美貌来装饰自己，而是以永恒和宝贵的恩宠。我们相信，凡是逃避那些似乎美善的事物，而紧跟那些被认为困难和奇怪的事物的人，就进入了真福。因为\OriginalTerm{恶行}{Vice}，当她通过模仿那不朽坏的（因为真正不朽坏的，她既没有也不能产生）时，用\OriginalTerm{德行}{Virtue}的属性和真正卓越的质素作为伪装包裹她自己的行为，就掳掠了属世心思的人，把自己的邪恶属性附加给\OriginalTerm{德行}{Virtue}。但那些理解属于真实之物的卓越性的人，在德行上也是不朽坏的。每一个有理性的人都应当这样想：关于基督徒、竞技运动员，以及那些做了诗人们所述所谓神的事的人，从我们对死亡的藐视——即使死亡可以逃脱——得出同样的结论。\par

% structure: p0025 source=h2 target=section reason=relative-heading-rank; base=h2

\section{第十二章 基督徒因藐视死亡而被证明无罪}

因为我自己，当我正陶醉于\Person{柏拉图}的学说时，听到基督徒被诽谤，又看到他们不畏死亡和一切被认为可怕的事物，就察觉到他们不可能生活在邪恶和享乐中。因为哪个好色或放纵的人，或哪个认为以人肉为食是好的，会欢迎死亡从而被剥夺其享乐，而不愿意继续当前的生活，并试图逃避统治者的注意呢？更不用说，当后果是死亡时，他会指控自己了。这也是恶鬼现在藉着恶人促使发生的事情。因为有些人因对我们的虚假控告而被处死，他们还把我们的家人，无论是孩子还是虚弱的妇女，拖去受刑，并用可怕的酷刑强迫他们承认那些他们自己公开犯下的荒诞行为；我们对此并不太在意，因为这些行为没有一件真是我们的，并且我们有那无始无终的、不可言说的天主作我们思想和行为的见证。因为我们为什么不甚至公开宣称这些是我们认为好的事，并证明这些是神圣的哲学，说\Person{萨图恩}的秘仪是在我们杀一个人时执行的，并且当我们像据说那样喝饱了血时，我们是在做你们在你们敬拜的偶像面前所做的事，你们不仅在那偶像上洒了无理性的动物的血，也洒了人的血，并用你们中间最杰出高贵的人的手所杀之人的血作奠祭呢？并且效仿\Person{朱庇特}和其他神祇在鸡奸和与女人无耻交合的事上，我们难道不能以\Person{伊壁鸠鲁}和诗人们的作品作为我们的辩护吗？但因为劝勉人们避开这样的教导，以及所有实行这些事并效法这些榜样的人，正如在这篇论述中我们努力劝服你们一样，我们受到各种方式的攻击。但我们不担忧，因为我们知道天主是万有的公正观察者。但愿现在有人能登上高高的讲坛，大声呼喊：\Quote{你们羞愧吧，羞愧吧，你们这些指控无辜者犯有你们自己公开可能犯下的罪行的人，并把适用于你们自己和你们的神祇的事归给那些对这些事没有丝毫兴趣的人。皈依吧，变聪明吧。}\par

% structure: p0027 source=h2 target=section reason=relative-heading-rank; base=h2

\section{第十三章 圣言如何曾在众人之内}

因为我自己，当我发现恶神围绕基督徒的神圣教义所抛出的邪恶伪装，为的是使他人远离加入他们时，我就嘲笑那些编造这些谎言的人，也嘲笑伪装本身和大众意见；并且我承认，我既夸耀，也竭尽全力要被人视为基督徒；不是因为柏拉图的教训与基督的教训不同，而是因为它们并非在所有方面都相似，正如其他教训——斯多葛派、诗人和历史学家的教训——也不完全相似。因为每个人按照自己分得种子圣言的比例，看到了与之相关的事物，从而说得很好。但是那些在更重要的问题上自相矛盾的人，似乎并未拥有属天的智慧，和那无可辩驳的知识。凡在众人中正确说出的，都是我们基督徒的财产。因为在天主之后，我们朝拜并爱慕那出自无始和不可言喻的天主的圣言，因为祂也为了我们而成为人，好使祂分担我们的痛苦，也带给我们医治。因为所有的作者，借着在他们内植入的圣言的播种，都能隐约看见现实。因为按照能力所接受的种子和模仿是一回事，而事物本身——即借着来自祂的恩宠而有的分沾和模仿——则是全然另一回事。\par

% structure: p0029 source=h2 target=section reason=relative-heading-rank; base=h2

\section{第十四章。儒斯定祈求将此陈情书公布}

因此我们恳求你公布这本小书，附上你认为合适的内容，好让我们的意见为他人所知，并使那些因自身过错而应受惩罚的人，能有公平的机会从错误的观念和对善的无知中获得释放；这样，这些事便可昭示于人，因为认识善恶是人的本性；而他们谴责他们不了解的我们，因为我们做他们说邪恶的事，并因喜欢那些做了此类事情的神祇——这些神祇甚至如今还要求人做类似的事——并且因为他们加给我们死亡、锁链或其他类似惩罚，仿佛我们犯了这些罪，他们便定了自己的罪，因此不需要别的法官了。\par

% structure: p0031 source=h2 target=section reason=relative-heading-rank; base=h2

\section{第十五章。结论}

并且我鄙视了我本族西蒙的邪恶而欺骗性的教训。如果你赋予此书权威，我们将在众人面前揭露他，以便他们如果能的话可以悔改。我们撰写这篇论述的唯一目的即在于此。而我们的教训，根据冷静的判断，并非可耻，反而比所有人类哲学更高超；即使不是如此，至少也与索塔德主义者、菲莱尼德信徒、舞蹈者（教派）、伊壁鸠鲁派以及诗人们的其他这类教训不同——所有人都被允许同时以表演和文字的方式熟悉这些教训。从此我们将保持沉默，因为我们尽了自己所能，并附加了祈祷：愿所有地方的人都堪当蒙受真理。并且我们也愿你，以合乎虔敬与哲学的方式，为你自己的缘故公正判断！\par

\SourceRecord{Translated by Marcus Dods and George Reith.；Ante-Nicene Fathers；Vol. 1.；Edited by Alexander Roberts, James Donaldson, and A. Cleveland Coxe.；Buffalo, NY: Christian Literature Publishing Co.,；1885.；<http://www.newadvent.org/fathers/0127.htm>.}
